% ======================================================================
% Chapter 14: Invitation: The Parton Model of Hadron Structure
% (book pages 473-480; PDF pages 496-503)
% ======================================================================
\chapter{引子：强子结构的部分子模型}

在本书第~II 部分中，我们以形式化的方式探索了量子场论的结构。我们发展了精致的计算方法（第~10 章），导出了提取标度律与渐近行为的形式体系（第~12 章），并推导了自发对称性破缺的一些推论（第~11 章）。这些形式体系中有相当多内容被证明在统计力学中有着出人意料的应用。然而，我们尚未研究它对基本粒子物理学的意义。为此，我们必须首先问：究竟是哪些特定的量子场论描述了基本粒子的相互作用。

自 20 世纪 70 年代中期以来，大多数高能物理学家都同意，构成物质的基本粒子是一组费米子，它们主要通过交换矢量玻色子而相互作用。基本费米子包括轻子（电子、其较重的对应粒子 $\mu$ 和 $\tau$，以及与上述每一种粒子对应的一个中性、几乎无质量的中微子），以及夸克，夸克的束缚态构成具有核相互作用的粒子——介子和重子（统称为强子）。这些费米子通过三种力相互作用：强相互作用、弱相互作用和电磁相互作用。其中，强相互作用负责核束缚以及原子核各成分之间的相互作用，而弱相互作用负责放射性 $\beta$ 衰变过程。电磁相互作用就是我们所熟悉的量子电动力学，它最小耦合到所有带电夸克和轻子。这三种力是否足以解释基本费米子最微妙的性质尚不清楚——我们将在第~20 章讨论这一问题——但这三种力当然是最突出的。现在人们认识到，这三种力都由矢量玻色子的交换来传递。描述电磁相互作用的方程由 Maxwell 发现，其量子力学含义已在第~I 部分中详细处理。弱相互作用和强相互作用的正确理论发现得要晚得多。

到 20 世纪 50 年代末，对弱相互作用截面与衰变率中螺旋度依赖的研究已经表明，弱相互作用涉及由夸克场和轻子场构成的矢量流之间的耦合。%
\footnote{关于弱相互作用唯象学的综述，参见 Perkins (1987) 第~7 章，或任何其他现代粒子物理教科书。} 因此很自然地假设弱相互作用是由非常重的矢量玻色子的交换引起的；事实上，这样的玻色子——$W$ 和 $Z$ 粒子——已于 1982 年在 CERN 的实验中被发现。但是，一个完整的弱相互作用理论不仅必须包含玻色子与费米子之间正确的耦合，还必须包含玻色子场自身的运动方程，即对 $W$ 和 $Z$ 而言的 Maxwell 方程组的类比物。要找到这些方程的正确形式并非易事，因为 Maxwell 方程组不允许矢量粒子获得质量。要正确调和广义 Maxwell 方程组与不为零的 $W$、$Z$ 质量，结果要求把自发对称性破缺纳入该理论。第~20 章和第~21 章将相当详细地处理这一课题，描述矢量场论与自发对称性破缺之间的相互作用。这种相互作用导致新的曲折和新的现象，超出了我们在第~11 章处理自发对称性破缺时所讨论的范围。一个完整的弱相互作用理论还要求同时纳入电磁相互作用，形成一种统一的结构，正如 Glashow、Weinberg 和 Salam 首先假设的那样。

另一方面，交换矢量玻色子的理论能否正确描述强相互作用，在很长一段时间里完全是模糊不清的。谜团的一部分在于，夸克并不以孤立的粒子形式存在。它们的存在，以及最终它们的量子数，都必须从可观测的强相互作用粒子谱中推断出来。但此外，还存在由强相互作用很强这一事实引起的复杂性。费曼图展开假设耦合常数很小；当耦合变强时，大量图都会变得重要（如果该级数确实收敛的话），从而不可能挑出基本相互作用顶角的贡献。强相互作用具有矢量性的关键线索，来自一个起初看来似乎只是另一个谜团的现象，即强相互作用在动量转移很大时会自行关闭，其含义我们现在就来描述。

\subsection*{几乎自由的部分子}

在第~5.1 节中，我们计算了 QED 过程 $e^{+}e^{-}\to\mu^{+}\mu^{-}$ 的截面。然后我们指出，$e^{+}e^{-}$ 湮灭为强子的对应截面也可以用同样的方法计算，只需采用一个把夸克当作无相互作用费米子处理的简单模型。这个方法给出了一个惊人准确的截面公式，抓住了其最重要的定性特征。但我们推迟了对这一谜团的解释：一个无相互作用夸克的模型，如何能代表一种在其他情形下极其强大的力的行为呢？

事实上，在高能强相互作用的研究中，有许多情形这种力的效应都出乎意料地弱。历史上，第一个这样的情形出现在质子-质子碰撞中。在高能下，当质心系能量高于 10 GeV 左右时，质子（或任何其他强子）的碰撞会产生大量 $\pi$ 介子。人们或许会想象这些 $\pi$ 介子会填满所有允许的相空间，但事实上，它们主要产生于动量几乎与碰撞轴共线的情形。产生一个具有大的横向于碰撞轴的动量分量的 $\pi$ 介子的概率，随该横向动量的值指数衰减，从而对于大于几百 MeV 的横向动量大大抑制了这种产生。

这种\emph{有限横向动量}现象导致了一种图像，把强子看成由许多组分松散束缚而成的集合。在这种图像中，被另一个质子撞击的质子会被撕碎成一团碎片。这些碎片的动量大致与质子的原始动量共线，并最终重新形成沿碰撞轴运动的强子。按假设，这些碎片不能吸收大的动量转移。我们可以从数学上这样刻画这一假设：在高能碰撞中，两个初始强子的动量几乎是类光的。沿碰撞轴排列的强子碎片也具有与原始动量矢量平行的类光动量。这种末态可以通过在碎片之间交换动量 $q$ 来产生，方式使得尽管 $q$ 的分量可能很大，不变量 $q^2$ 却总是很小。以大横向动量抛出一个强子需要大的（类空的）$q^2$，但这样的过程非常罕见。因此人们假设，强子是像果冻一样松散的组分云，不能吸收大的 $q^2$。

20 世纪 60 年代末，SLAC-MIT 深度非弹性散射实验对这种强子结构图像进行了关键检验。%
\footnote{关于这些实验及其影响的描述，参见 J.~I.~Friedman, H.~W.~Kendall 和 R.~E.~Taylor, \emph{Rev.\ Mod.\ Phys.} \textbf{63}, 573 (1991)。} 在这些实验中，一束 20 GeV 的电子束被氢靶散射，并对大的偏转角测量散射率，这对应于从电子到靶中质子的大不变量动量转移。大的动量转移是通过电磁相互作用而非强相互作用传递的，因此所传递的动量大小可以由被散射电子的动量计算出来。在把强子视为复杂且弱束缚的模型中，人们预期散射率会非常低。

然而，SLAC-MIT 实验却观察到电子在质子上硬散射的显著速率。总反应率与假如质子是一个按 QED 最简单预期散射的基本粒子时所预期的大小相当。然而，只有在罕见的情形下，散射过程中才会有一个单独的质子出现。速率的最大部分来自相空间的\emph{深度非弹性}区域，其中电磁冲量把质子打碎，并产生一个含有大量强子的系统。

人们怎样才能调和电磁硬散射过程的存在与强相互作用过程中几乎完全没有硬散射这一事实呢？为了回答这个问题，Bjorken 和 Feynman 提出了如下简单模型，称为\emph{部分子模型}：假设质子是由少数组分松散束缚而成的集合，这些组分称为\emph{部分子}。其中包括夸克（和反夸克），它们是一类携带电荷的费米子，还可能包括其他负责把它们束缚在一起的中性粒子。按假设，这些组分不能通过强相互作用交换大的动量 $q^2$。然而，夸克具有基本费米子的电磁相互作用，因此从夸克上散射的电子可以把夸克撞出质子。被撞的夸克随后与质子的其余部分软交换动量，使得质子的碎片以强子喷注的形式物质化。所产生的强子应与原始被撞部分子的方向共线。

部分子模型尽管不完整，却对深度非弹性电子散射的截面施加了一个很强的约束。为了导出这一约束，我们首先考虑电子在单个组分夸克上散射的截面。我们在第~5.4 节中讨论过电子-$\mu$ 子散射这一相关过程，可以借用那里的结果。由于我们设想反应在非常高的能量下发生，我们将忽略所有质量。在无质量极限下，不变矩阵元平方以简洁的形式写在式~(5.71) 中：
\begin{equation}
  \frac{1}{4}\sum_{\text{spins}}|\mathcal{M}|^2
  = \frac{4\pi^2\alpha^2 Q_i^2}{s^2}\ \frac{s^2+u^2}{t^2} ,
  \label{eq:14.1}
\end{equation}
其中 $s$、$t$、$u$ 是电子-夸克碰撞的 Mandelstam 变量，$Q_i$ 是以 $|e|$ 为单位的夸克电荷。回顾式~\eqref{eq:5.73}，对于涉及无质量粒子的碰撞，$s + t + u = 0$。于是质心系中的微分散射截面为
\begin{equation}
  \frac{d\sigma}{dt} = \frac{2\pi\alpha^2 Q_i^2}{s^2}\ \frac{s^2+u^2}{t^2}
  .
  \label{eq:14.2}
\end{equation}
或者，由于 $t = -s(1-\cos\theta_{cm})/2$，
\begin{equation}
  \frac{d\sigma}{d\Omega_{cm}} = \frac{\alpha^2 Q_i^2}{4s}\ \frac{1 +
    \cos^2\theta_{cm}}{\sin^4(\theta_{cm}/2)}
  \ \cdots .
  \label{eq:14.3}
\end{equation}

\begin{figure}[t]
  \centering
  \includegraphics[width=0.55\textwidth]{fig14.1.png}
  \caption{部分子模型中深度非弹性电子散射的运动学。}
  \label{fig:14.1}
\end{figure}

为了利用这一结果，我们必须把不变量 $s$ 和 $t$ 与电子-质子非弹性散射的实验可观测量联系起来。运动学变量如图~\ref{fig:14.1} 所示。来自电子的动量转移 $q$ 可以通过测量电子的末态动量和能量来测定，无需使用强子产物的任何信息。由于 $q^\mu$ 是一个类空矢量，通常用正量 $Q$ 来表示它的不变量平方，有
\begin{equation}
  Q^2 = -q^2 .
  \label{eq:14.4}
\end{equation}
于是不变量 $t$ 就是 $-Q^2$。

用可测量量来表示 $s$ 则更困难。如果从电子-质子质心系来观察碰撞，并把质子想象成由部分子松散束缚而成的集合（并继续忽略质量），我们可以用某个部分子所携带的质子总动量份额来刻画它。我们用参数 $\xi$ 表示这一纵向份额，其中 $0 < \xi < 1$。对每一种类型的部分子 $i$，例如电荷 $Q_i = +2/3$ 的 up 型夸克，都有一个函数 $f_i(\xi)$，表示质子含有一个类型为 $i$、纵向份额为 $\xi$ 的部分子的概率。电子-质子非弹性散射总截面的表达式将包含对被打部分子 $\xi$ 值的积分。于是部分子的动量矢量为 $p = \xi P$，其中 $P$ 是质子的总动量。因此，若 $k$ 是电子的初态动量，
\begin{equation}
  s = (p+k)^2 = 2p\cdot k = 2\xi P\cdot k = \xi s ,
  \label{eq:14.5}
\end{equation}
其中 $s$ 是电子-质子质心系能量的平方。

引人注目的是，只要假设电子-部分子散射是弹性的，$\xi$ 也可以仅通过测量电子动量来确定。由于被散射的部分子的质量与 $s$ 和 $Q^2$ 相比很小，
\begin{equation}
  0 \approx (p+q)^2 = 2p\cdot q + q^2 = 2\xi P\cdot q - Q^2 .
  \label{eq:14.6}
\end{equation}
因此
\begin{equation}
  \xi = x,
  \qquad\text{其中}\qquad
  x = \frac{Q^2}{2P\cdot q} .
  \label{eq:14.7}
\end{equation}
由每一个被散射的电子，人们可以确定该散射过程的 $Q^2$ 和 $x$ 的值。部分子模型于是预言 $x$--$Q^2$ 平面上的事例分布。利用在 $\xi = x$ 处取值的部分子分布函数 $f_i(\xi)$ 以及截面公式 \eqref{eq:14.3}，我们得到分布
\begin{equation}
  \frac{d^2\sigma}{dx\,dQ^2}
  = \frac{2\pi\alpha^2}{Q^4}\ \sum_i Q_i^2\,f_i(x)\ \Bigl[ 1 + \Bigl(
    1 - \frac{Q^2}{xs} \Bigr)^2 \Bigr] .
  \label{eq:14.8}
\end{equation}
分布函数 $f_i(x)$ 依赖于质子结构的细节，人们尚不知道如何从第一性原理计算它们。但公式 \eqref{eq:14.8} 仍然作出了一个惊人的预言，即深度非弹性散射截面在除以因子
\begin{equation}
  \frac{2\pi\alpha^2}{Q^4}\ \Bigl[ 1 + \Bigl( 1 - \frac{Q^2}{xs}
    \Bigr)^2 \Bigr]
  \label{eq:14.9}
\end{equation}
以去除 QED 截面的运动学依赖之后，给出的量只依赖于 $x$ 而与 $Q^2$ 无关。这种行为称为\emph{Bjorken 标度无关性}。事实上，SLAC-MIT 实验的数据在 $Q$ 值高于 1 GeV 时以约 10\% 的精度表现出 Bjorken 标度无关性，如图~\ref{fig:14.2} 所示。

\begin{figure}[t]
  \centering
  \includegraphics[width=0.6\textwidth]{fig14.2.png}
  \caption{利用 SLAC-MIT 实验测得的 $e^{-}p$ 深度非弹性散射截面检验 Bjorken 标度无关性，J.~S.~Poucher 等人，\emph{Phys.\ Rev.\ Lett.} \textbf{32}, 118 (1974)。针对所标出的各种初态电子能量和散射角，我们把 $d^2\sigma/dx\,dQ^2$ 除以因子 \eqref{eq:14.9} 后对 $x$ 作图。数据覆盖 $1\ \mathrm{GeV}^2 < Q^2 < 8\ \mathrm{GeV}^2$ 的范围。}
  \label{fig:14.2}
\end{figure}

Bjorken 标度无关性本质上是在说，无论质子被撞得多狠，对电磁探针而言质子的结构看起来都相同。在质子的静止系中，交换的虚光子的能量为
\begin{equation}
  q^0 \approx \frac{1}{2m}\bigl( Q^2 + \cdots \bigr) ,
\end{equation}
其中 $m$ 是质子的质量。该能量转移的倒数大致就是质子各组分所看到的散射过程的持续时间。这个时间应当与质子质量的倒数相比较，后者是部分子相互作用的特征时间。深度非弹性区发生在 $q^0 \gg m$ 时，即散射相对于质子的正常时间尺度非常迅速时。Bjorken 标度无关性意味着，在这样一次迅速的散射过程中，质子各组分之间的相互作用可以忽略。我们可以想象，在对应于 1 GeV 或更高能量转移的极短时间尺度上，部分子近似为自由粒子，尽管在更长的时间尺度上它们有强相互作用。

\subsection*{渐近自由的部分子}

Bjorken 标度无关性所隐含的质子结构图像简单而优美，但它提出了新的、根本性的问题。在量子场论中，费米子通过交换虚粒子相互作用。这些虚粒子可以
\begin{equation}
  \text{(在短距离处增大耦合的有效强度)} .
\end{equation}
但是，对深度非弹性散射数据的分析却表明，强相互作用在短距离处会变\emph{弱}。20 世纪 70 年代初，这一表面上的矛盾被 Gross、Politzer 和 Wilczek 发现的\emph{渐近自由}所解决。他们证明，非阿贝尔规范理论——即相互作用的矢量玻色子的理论——具有一个显著的性质，即有效耦合常数在大动量处趋于零。这正是在四维强相互作用的量子场论中使 Bjorken 标度无关性有意义所要求的行为。这一发现为构建强相互作用的基本理论提供了关键线索。显然，夸克正是被这一类型的相互作用的矢量玻色子（称为\emph{胶子}）束缚在一起的。

然而，这些规范理论并不能精确再现严格 Bjorken 标度无关性的预期。自由部分子模型与渐近自由的量子场论模型之间的差异，出现在人们把深度非弹性散射以及其他涉及大动量转移的强相互作用过程的测量精度提高一个层次时。在渐近自由的量子场论中，耦合常数在任何有限的动量转移处仍然不为零。事实上，耦合常数向零的最终演化非常缓慢，是对动量呈对数式的。因此，在某个层次上，人们必然会发现对 Bjorken 标度无关性的小修正，它们与高动量胶子的交换或发射有关。类似地，高动量转移下强子物理学的其他定性简化——例如，强子-强子碰撞中有限横向动量的现象——应当只是近似的，会收到来自胶子交换与发射的修正。因此，强相互作用渐近自由理论的预言有两方面。一方面，这样的理论预言高动量下行为的定性简化。但另一方面，这样的理论又预言了这种行为的一个特定的修正模式。

事实上，20 世纪 70 年代的粒子物理实验恰恰揭示了这样一幅图像。人们发现 Bjorken 标度无关性只是一个近似关系，表现出对应于部分子分布 $f_i(x)$ 在 $Q^2$ 的对数标度上缓慢演化的破坏。人们发现，强子-强子碰撞中粒子产生的速率在横向动量非常大时只按幂律而非指数方式下降，并且大横向动量处产生的粒子被证明与硬散射夸克或胶子软演化所产生的强子喷注相关。最引人注目的是，人们为这些以及其他偏离标度无关的现象所找到的截面形式，最终为传递强相互作用的基本场的矢量性提供了直接证据。

我们将在第~17 章回顾所有这些现象，届时我们将研究描述强相互作用的特定规范理论。然而，首先我们必须学会如何构造非阿贝尔规范理论，以及如何用费曼图算出它们的预言。在我们对这些理论的全部分析中，重整化群将扮演至关重要的角色。研究非阿贝尔规范理论的一个非常优美之处在于，量子场论最有力的一般思想在与这些结构精巧的特定模型的特殊性质交织在一起时，会获得更大的力量。一般原理与规范理论特殊性质之间的这种相互作用，将是本书第~III 部分的主要主题。
